\documentclass[12pt,a4paper]{article}
\usepackage[utf8]{inputenc}
\usepackage[T1]{fontenc}
\usepackage[brazil]{babel}
\usepackage[left=3cm,right=2cm,top=3cm,bottom=2cm]{geometry}
\usepackage{fancyhdr}
\usepackage{setspace}
\usepackage{parskip}

% ==========================================
% CABEÇALHO E RODAPÉ
% ==========================================
\pagestyle{fancy}
\fancyhf{}
\fancyhead[C]{%
  \textbf{\large AUTORIDADE PORTUÁRIA}\\
  \small Assessoria Jurídica
}
\fancyfoot[C]{\small Página \thepage}
\fancyfoot[L]{\small Documento gerado pelo PortoGpt}
\renewcommand{\headrulewidth}{0.4pt}
\renewcommand{\footrulewidth}{0.4pt}

% ==========================================
% INÍCIO DO DOCUMENTO
% ==========================================
\begin{document}
\onehalfspacing

\begin{center}
  \vspace*{1cm}
  {\LARGE \textbf{PARECER TÉCNICO-JURÍDICO}}\\[0.5cm]
  {\large Nº 001/2026 -- ASJUR}\\[1cm]
\end{center}

% --- Dados do Parecer ---
\begin{tabular}{@{}ll@{}}
  \textbf{Processo:} & Proc. Adm. nº 2026.001.00123 \\
  \textbf{Interessado:} & Diretoria de Administração e Finanças \\
  \textbf{Assunto:} & Legalidade da Contratação Emergencial de Serviços \\
  \textbf{Data:} & 07 de maio de 2026 \\
\end{tabular}

\vspace{1cm}
\hrule
\vspace{0.5cm}

\section*{I -- DO RELATÓRIO}

Trata-se de consulta formulada pela Diretoria de Administração e Finanças acerca da possibilidade jurídica de contratação emergencial de serviços de manutenção preventiva nos sistemas de combate a incêndio do Terminal Portuário.

A consulta é motivada pela constatação, em inspeção recente realizada pelo Corpo de Bombeiros Militar, de que o sistema hidráulico de combate a incêndio apresenta falhas que comprometem a segurança das operações portuárias.

\section*{II -- DA FUNDAMENTAÇÃO}

A contratação emergencial encontra respaldo no art.\ 75, inciso VIII, da Lei nº 14.133/2021 (Nova Lei de Licitações), que autoriza a dispensa de licitação nos casos de emergência ou calamidade pública que possam ocasionar prejuízo ou comprometer a continuidade dos serviços públicos ou a segurança de pessoas.

Para a configuração da emergência, devem estar presentes, cumulativamente, os seguintes requisitos:

\begin{enumerate}
  \item Situação de emergência devidamente caracterizada e justificada;
  \item Risco de prejuízo ou comprometimento da segurança;
  \item Impossibilidade de aguardar o trâmite regular do procedimento licitatório;
  \item Contratação pelo prazo máximo de 1 (um) ano.
\end{enumerate}

\section*{III -- DA CONCLUSÃO}

Diante do exposto, esta Assessoria Jurídica opina pela \textbf{viabilidade jurídica} da contratação emergencial pretendida, desde que observadas as seguintes condições:

\begin{itemize}
  \item Elaboração de justificativa técnica detalhada pela área competente;
  \item Realização de pesquisa de preços com, no mínimo, três orçamentos;
  \item Limitação do prazo contratual a 180 (cento e oitenta) dias;
  \item Publicação do extrato do contrato no Diário Oficial.
\end{itemize}

É o parecer, salvo melhor juízo.

\vspace{1.5cm}

\begin{center}
\textbf{Dr.\ Carlos Mendes}\\
OAB/MA 54321\\
Assessor Jurídico da Autoridade Portuária
\end{center}

\end{document}
